\section{Problem 786}

\subsection*{1) ROUND-2 OBJECTIVE}
Path \textbf{(C)}: identify the fatal obstruction in Round 1 and correct the statement. In Round 1 I treated the repetitions-allowed version as still open. That is incorrect. The currently curated literature record shows that, when repetitions are allowed, the problem is already settled negatively: every such set is contained in a level set of a real-valued additive function, and Erd\H{o}s--Ruzsa--S\'ark\"ozy imply an absolute upper-density barrier. The minimal corrected open problem is the \emph{distinct-elements} version.

\subsection*{2) Round-1 FOUNDATION USED}
I use the following Round-1 results without reproving them.
\begin{enumerate}
\item The finite additive-functional criterion: for finite $A$, the repetitions-allowed property is equivalent to the existence of a rational additive functional on prime-exponent vectors taking the value $1$ on every $a \in A$.
\item The Round-1 diagnosis that the real gap was to control level sets of additive functions.
\item The standard examples (the class $n \equiv 2 \pmod 4$, Selfridge's construction, and the one-large-prime finite construction).
\end{enumerate}

\subsection*{3) NEW INSIGHT / TOOL (ROUND-2)}
The new ingredient is an \emph{infinite-set extension} of the Round-1 structural lemma. Instead of working only with finitely many primes and rational coefficients, I pass to the rational vector space of finitely supported prime-exponent vectors and extend a linear functional from the span of $\{\mathbf v(a): a \in A\}$ to the whole space. This gives a real-valued additive function on all natural numbers.

The second new ingredient is the corrected literature input: the updated 2026 discussion of Problem 786 records that Erd\H{o}s--Ruzsa--S\'ark\"ozy already settle the repetitions version negatively, and that the remaining open variant is the distinct-elements interpretation.

\subsection*{4) ATTACK PLAN (ROUND-2)}
The exact gap left by Round 1 was this: the finite linear-algebra lemma had not yet been upgraded to arbitrary $A$, so it could not be combined with the known distribution theory of additive functions. I will:
\begin{enumerate}
\item extend the structural lemma from finite $A$ to arbitrary $A$;
\item reduce the repetitions-allowed problem to a level set $\{n : f(n)=1\}$ of a real additive function $f$;
\item import the Erd\H{o}s--Ruzsa--S\'ark\"ozy consequence recorded in the current literature summary, obtaining a negative answer to both the density and finite problems;
\item isolate precisely why this argument does \emph{not} resolve the distinct-elements variant.
\end{enumerate}

\subsection*{5) WORK (ROUND-2)}
Write the repetitions-allowed property as follows:
\[
 a_1 \cdots a_r = b_1 \cdots b_s \quad (a_i,b_j \in A)
 \implies r=s.
\]

\paragraph{Lemma 786.1 (global additive-functional criterion).}
Let $A$ be a subset of the natural numbers with the repetitions-allowed property. Then there exists a real-valued additive function $f$ on the natural numbers such that
\[
 f(a)=1 \qquad \text{for every } a \in A.
\]

\paragraph{Proof.}
Let $P$ be the set of primes dividing at least one element of $A$. Let $V$ be the vector space over the rationals consisting of all finitely supported tuples
\[
 x=(x_p)_{p \in P}, \qquad x_p \in \mathbf Q.
\]
For a natural number $n$, let
\[
 \mathbf v(n) := (v_p(n))_{p \in P} \in V,
\]
where $v_p(n)$ is the exponent of $p$ in the factorisation of $n$.

Let $W$ be the rational span of $\{\mathbf v(a): a \in A\}$ inside $V$. Define a map $L$ on $W$ by declaring
\[
 L\!\left(\sum_{i=1}^m q_i \mathbf v(a_i)\right) := \sum_{i=1}^m q_i
\]
for rational numbers $q_i$ and elements $a_i \in A$.

I claim this is well-defined. Suppose
\[
 \sum_{i=1}^m q_i \mathbf v(a_i)=0.
\]
Choose a common denominator $D$ and set $m_i = D q_i \in \mathbf Z$. Then
\[
 \sum_{i=1}^m m_i \mathbf v(a_i)=0.
\]
Write $m_i = m_i^+ - m_i^-$ with $m_i^+,m_i^- \ge 0$. The equality of exponent vectors is equivalent to
\[
 \prod_{i=1}^m a_i^{m_i^+} = \prod_{i=1}^m a_i^{m_i^-}.
\]
Each side is a product of elements of $A$, with repetitions allowed. By the defining property of $A$, the total number of factors on the two sides must be equal:
\[
 \sum_{i=1}^m m_i^+ = \sum_{i=1}^m m_i^-.
\]
Hence
\[
 \sum_{i=1}^m m_i = 0,
\]
and dividing by $D$ gives $\sum_i q_i = 0$. So $L$ is well-defined.

Now extend $L$ from the subspace $W$ to a linear functional on all of $V$ by choosing a basis of $W$ and extending it to a basis of $V$. Keep the same name $L$ for the extension. Define
\[
 f(n) := L(\mathbf v(n)).
\]
Because
\[
 \mathbf v(mn)=\mathbf v(m)+\mathbf v(n),
\]
we have
\[
 f(mn)=f(m)+f(n)
\]
for all natural numbers $m,n$; that is, $f$ is additive. And for every $a \in A$,
\[
 f(a)=L(\mathbf v(a))=1.
\]
This proves the lemma.
\hfill$\square$

\paragraph{Corollary 786.2.}
Every repetitions-allowed set $A$ is contained in a level set of a real-valued additive function:
\[
 A \subseteq \{n : f(n)=1\}
\]
for some additive $f$.

\paragraph{External input.}
The current 2026 update to the Erd\H{o}s Problems page for Problem 786 records the following consequence of Erd\H{o}s--Ruzsa--S\'ark\"ozy (1973): once a set is contained in a fixed level set of a real-valued additive function, one obtains
\[
 \overline d(A) \le \frac12,
\]
and, in the finite problem, there is an absolute constant $c>0$ such that for all sufficiently large $N$,
\[
 |A \cap [1,N]| \le (1-c)N.
\]
Applying this to Corollary 786.2 gives the promised negative answers.

\paragraph{Theorem 786.3 (repetitions version is already settled).}
For the repetitions-allowed interpretation of Problem 786:
\begin{enumerate}
\item there is no set $A$ of density $>1-\varepsilon$ for every $\varepsilon>0$; in fact every such $A$ has upper density at most $1/2$;
\item in the finite problem one cannot have $|A|=(1-o(1))N$; indeed $|A| \le (1-c)N$ for some absolute $c>0$ and all sufficiently large $N$.
\end{enumerate}

\paragraph{Why the distinct-elements version is genuinely different.}
The proof above uses arbitrary integer relations among exponent vectors. Such a relation becomes a product identity only because repetitions are allowed. If one insists that all $a_i$ and all $b_j$ be distinct, then only disjoint $0$--$1$ relations are relevant; integer coefficients larger than $1$ are no longer legal. So the additive-function reduction collapses exactly at this point. This isolates the minimal corrected open statement: require the products on both sides to use distinct elements.

\subsection*{6) ADVERSARIAL VERIFICATION}
\begin{enumerate}
\item The definition of $L$ uses only finitely many elements of $A$ at a time, so no infinite product or convergence issue appears.
\item The extension step is safe because $W$ is a rational vector subspace of $V$; I am extending a linear functional on vector spaces, not claiming that a homomorphism to the integers extends on a free abelian group. This avoids the standard obstruction exemplified by $2\mathbf Z \to \mathbf Z$.
\item The conclusion is only $A \subseteq \{f=1\}$, not equality; this is the correct statement needed for the upper-bound transfer.
\item The distinct-elements obstruction is real: the argument from an integer relation
\[
 \sum m_a \mathbf v(a)=0
\]
produces repeated factors $a^{m_a^+}$ and $a^{m_a^-}$, which are illegal in the distinct-elements variant whenever some coefficient has absolute value bigger than $1$.
\end{enumerate}

\subsection*{7) FINAL}
\textbf{FULL SOLUTION -- COUNTEREXAMPLE / DISPROOF.}

For the repetitions-allowed interpretation, both optimistic conclusions are false. Every such set is contained in a level set of a real-valued additive function, and the Erd\H{o}s--Ruzsa--S\'ark\"ozy consequences recorded in the current literature give upper density at most $1/2$ and a finite upper bound $|A| \le (1-c)N$. The minimal corrected open problem is the distinct-elements version.

\subsection*{8) COMPLETION ESTIMATE}
\textbf{COMPLETION: 100\%}

\subsection*{9) REFERENCES}
\begin{thebibliography}{99}
\bibitem{ERS73-786}
P. Erd\H{o}s, I. Ruzsa, and A. S\'ark\"ozy,
\emph{On the number of solutions of $f(n)=a$ for additive functions},
Acta Arith. 24 (1973), 1--9.

\bibitem{Bloom786}
T. F. Bloom,
\emph{Erd\H{o}s Problem \#786}, updated 2 February 2026,
\texttt{https://www.erdosproblems.com/786}.

\bibitem{Bloom786Thread}
T. F. Bloom, T. Tao, and others,
\emph{Discussion thread for Erd\H{o}s Problem \#786}, especially comments of 18 October 2025 and 2 February 2026,
\texttt{https://www.erdosproblems.com/forum/thread/786}.
\end{thebibliography}

